
\documentclass[11pt,a4paper]{article}

\usepackage{fontspec}
\usepackage[french]{babel}
\usepackage{amsmath,amssymb,amsthm,mathtools}
\usepackage{geometry}
\usepackage{xcolor}
\usepackage{fancyhdr}
\usepackage{enumitem}
\usepackage{hyperref}
\usepackage{titlesec}
\usepackage{microtype}
\usepackage{booktabs}
\usepackage{array}

\geometry{margin=2.2cm}
\setmainfont{DejaVu Serif}
\setsansfont{DejaVu Sans}
\setmonofont{DejaVu Sans Mono}

\definecolor{accent}{HTML}{1F4E79}
\definecolor{lightaccent}{HTML}{EAF2F8}

\pagestyle{fancy}
\fancyhf{}
\lhead{\small Cours Deep Reinforcement Learning}
\chead{\small Classe : 4IA}
\rhead{\small Année : 2026}
\lfoot{\small Enseignant : Mourad Zéraï, PhD - ESPRIT School of Engineering}
\cfoot{}
\rfoot{\thepage}
\setlength{\headheight}{14pt}

\titleformat{\section}{\large\bfseries\color{accent}}{\thesection.}{0.5em}{}
\titleformat{\subsection}{\normalsize\bfseries\color{accent}}{\thesubsection.}{0.5em}{}
\titleformat{\subsubsection}{\normalsize\bfseries}{\thesubsubsection.}{0.5em}{}

\newtheorem*{definition}{Définition}
\newtheorem*{rappel}{Rappel}
\newtheorem*{hypothese}{Hypothèse}
\newtheorem*{propriete}{Propriété}
\newtheorem*{theoreme}{Théorème}
\newtheorem*{remarque}{Remarque}

\newcommand{\E}{\mathbb{E}}
\newcommand{\Pbb}{\mathbb{P}}
\newcommand{\R}{\mathbb{R}}
\newcommand{\argmax}{\operatorname*{argmax}}
\newcommand{\given}{\,|\,}

\begin{document}

\begin{center}
    {\LARGE\bfseries Dérivation et explication pas à pas de Value Iteration}\\[0.4em]
    {\large Note de cours rigoureuse, détaillée et appliquée à FrozenLake}
\end{center}

\vspace{0.6em}

\section{Objectif}

Le but de cette note est d'expliquer, pas à pas, l'algorithme de \emph{Value Iteration}, de justifier rigoureusement chaque passage mathématique, puis de montrer comment l'appliquer à l'environnement FrozenLake.

\section{Cadre de départ}

On considère un MDP fini $(\mathcal S,\mathcal A,p,r,\gamma)$ avec :
\begin{itemize}[leftmargin=1.6em]
    \item un ensemble fini d'états $\mathcal S$ ;
    \item un ensemble fini d'actions $\mathcal A$ ;
    \item une dynamique $p(s',r \given s,a)$ ;
    \item un facteur d'actualisation $\gamma \in [0,1)$.
\end{itemize}

\begin{definition}[Valeur optimale]
La fonction de valeur optimale est
\[
v_*(s) = \sup_{\pi} v_{\pi}(s),
\]
où $v_{\pi}(s)$ désigne la valeur de l'état $s$ sous la politique $\pi$.
\end{definition}

\begin{definition}[Équation d'optimalité de Bellman]
Pour tout état $s$,
\[
v_*(s)=\max_{a\in\mathcal A}\sum_{s',r} p(s',r\given s,a)\left[r+\gamma v_*(s')\right].
\]
\end{definition}

L'algorithme de Value Iteration consiste à transformer cette équation en un schéma itératif.

\section{Construction de l'opérateur de Bellman optimal}

\subsection{Définition de l'opérateur}

On définit l'opérateur $T_*$ sur l'espace des fonctions $v:\mathcal S \to \R$ par
\[
(T_*v)(s)=\max_{a\in\mathcal A}\sum_{s',r} p(s',r\given s,a)\left[r+\gamma v(s')\right].
\]

\begin{rappel}
L'équation d'optimalité de Bellman s'écrit alors simplement
\[
v_* = T_* v_*.
\]
Autrement dit, $v_*$ est un point fixe de $T_*$.
\end{rappel}

\subsection{Pourquoi cette réécriture est utile}

Au lieu de chercher directement $v_*$, on part d'une fonction initiale $v^{(0)}$ puis on applique récursivement l'opérateur :
\[
v^{(k+1)} = T_* v^{(k)}.
\]
Il faut maintenant justifier pourquoi cette suite converge.

\section{Propriété de contraction}

\begin{theoreme}[Contraction de l'opérateur de Bellman optimal]
Pour la norme infinie
\[
\|v-w\|_{\infty} = \max_{s\in\mathcal S}|v(s)-w(s)|,
\]
on a
\[
\|T_*v - T_*w\|_{\infty} \le \gamma \|v-w\|_{\infty}.
\]
\end{theoreme}

\paragraph{Étape 1.}
Fixons un état $s$. Alors
\[
(T_*v)(s)=\max_a Q_v(s,a), \qquad (T_*w)(s)=\max_a Q_w(s,a),
\]
où
\[
Q_v(s,a)=\sum_{s',r} p(s',r\given s,a)\left[r+\gamma v(s')\right].
\]

\paragraph{Étape 2.}
On utilise l'inégalité élémentaire
\[
|\max_a x_a - \max_a y_a| \le \max_a |x_a-y_a|.
\]
Elle découle du fait que le maximum est une application 1-lipschitzienne pour la norme infinie sur un ensemble fini.

Donc
\[
|(T_*v)(s)-(T_*w)(s)| \le \max_a |Q_v(s,a)-Q_w(s,a)|.
\]

\paragraph{Étape 3.}
Pour une action fixée $a$,
\begin{align*}
Q_v(s,a)-Q_w(s,a)
&= \sum_{s',r} p(s',r\given s,a)\left[\gamma v(s')-\gamma w(s')\right] \\
&= \gamma \sum_{s',r} p(s',r\given s,a)\left[v(s')-w(s')\right].
\end{align*}
On a utilisé la linéarité de la somme et la distributivité du produit.

\paragraph{Étape 4.}
En valeur absolue,
\begin{align*}
|Q_v(s,a)-Q_w(s,a)|
&\le \gamma \sum_{s',r} p(s',r\given s,a)\, |v(s')-w(s')|.
\end{align*}
On a utilisé l'inégalité triangulaire.

\paragraph{Étape 5.}
Comme
\[
|v(s')-w(s')| \le \|v-w\|_{\infty} \quad \text{pour tout } s',
\]
on obtient
\begin{align*}
|Q_v(s,a)-Q_w(s,a)|
&\le \gamma \sum_{s',r} p(s',r\given s,a)\, \|v-w\|_{\infty}.
\end{align*}

\paragraph{Étape 6.}
Le terme $\|v-w\|_{\infty}$ ne dépend pas de $(s',r)$, donc
\[
|Q_v(s,a)-Q_w(s,a)| \le \gamma \|v-w\|_{\infty} \sum_{s',r} p(s',r\given s,a).
\]
Or, comme $p(\cdot,\cdot\given s,a)$ est une loi de probabilité,
\[
\sum_{s',r} p(s',r\given s,a)=1.
\]
Donc
\[
|Q_v(s,a)-Q_w(s,a)| \le \gamma \|v-w\|_{\infty}.
\]

\paragraph{Étape 7.}
En revenant à l'état $s$,
\[
|(T_*v)(s)-(T_*w)(s)| \le \gamma \|v-w\|_{\infty}.
\]
Comme ceci vaut pour tout état $s$, on prend le maximum sur $s$ et on obtient
\[
\|T_*v-T_*w\|_{\infty}\le \gamma \|v-w\|_{\infty}.
\]

\section{Conséquence : convergence de Value Iteration}

Comme $0\le \gamma <1$, l'opérateur $T_*$ est une contraction stricte sur l'espace vectoriel de dimension finie des fonctions sur $\mathcal S$, muni de la norme infinie. Le théorème du point fixe de Banach s'applique alors.

\begin{theoreme}[Point fixe et convergence]
Il existe un unique point fixe $v_*$ tel que
\[
v_* = T_* v_*,
\]
et pour toute initialisation $v^{(0)}$,
\[
v^{(k+1)} = T_*v^{(k)}
\]
converge vers $v_*$ lorsque $k\to\infty$.
\end{theoreme}

\section{Algorithme pas à pas}

\subsection{Initialisation}

On choisit une valeur initiale, par exemple
\[
v^{(0)}(s)=0 \quad \text{pour tout } s.
\]

\subsection{Mise à jour}

Pour chaque itération $k$ et pour chaque état $s$, on calcule
\[
v^{(k+1)}(s)=\max_{a}\sum_{s',r} p(s',r\given s,a)\left[r+\gamma v^{(k)}(s')\right].
\]

\subsection{Critère d'arrêt}

On calcule
\[
\Delta_k = \max_s |v^{(k+1)}(s)-v^{(k)}(s)|.
\]
On s'arrête lorsque $\Delta_k < \theta$ pour un seuil choisi.

\subsection{Extraction de la politique optimale}

Une fois la valeur stabilisée, on choisit pour chaque état
\[
\pi_*(s) \in \argmax_a \sum_{s',r} p(s',r\given s,a)\left[r+\gamma v_*(s')\right].
\]

\section{Application à FrozenLake}

\subsection{Description de l'environnement}

Dans FrozenLake 4x4 :
\begin{itemize}[leftmargin=1.6em]
    \item les états sont les 16 cases de la grille ;
    \item les actions sont Left, Down, Right, Up ;
    \item certaines cases sont des trous ;
    \item une case est le but ;
    \item la récompense vaut $1$ seulement quand l'agent atteint le but, et $0$ sinon ;
    \item les trous et le but sont terminaux.
\end{itemize}

Dans la version déterministe, l'action demandée est exécutée telle quelle. Dans la version glissante (\texttt{is\_slippery=True}), l'action choisie ne conduit pas toujours à la direction visée : la transition devient stochastique.

\subsection{Modèle utilisé}

FrozenLake fournit un modèle explicite des transitions de la forme
\[
P[s][a] = \{(p_i,s'_i,r_i,\text{done}_i)\}_i.
\]
Cela donne directement les probabilités $p(s',r\given s,a)$ nécessaires pour Value Iteration.

\subsection{Écriture concrète de la mise à jour}

Pour une case $s$ et une action $a$ :
\[
Q^{(k)}(s,a)=\sum_i p_i \left[r_i+\gamma v^{(k)}(s'_i)\right].
\]
Puis
\[
v^{(k+1)}(s)=\max_a Q^{(k)}(s,a).
\]

\subsection{Cas des états terminaux}

Si $s$ est terminal, alors il est naturel de fixer
\[
v^{(k+1)}(s)=0
\]
si l'on considère qu'aucune récompense supplémentaire n'est reçue après la terminaison.

\subsection{Interprétation}

Value Iteration propage progressivement la promesse du but vers les cases en amont. Dans FrozenLake glissant, la présence d'aléa réduit certaines valeurs car même une bonne action peut mener à un trou ou à une déviation.

\section{Pseudo-code pédagogique}

\begin{verbatim}
Initialiser v(s)=0 pour tous les états
Répéter :
    delta = 0
    pour chaque état s :
        old = v(s)
        v(s) = max_a somme_{s',r} p(s',r|s,a)[r + gamma v(s')]
        delta = max(delta, |v(s)-old|)
Jusqu'à ce que delta < theta

Pour chaque état s :
    pi*(s) = argmax_a somme_{s',r} p(s',r|s,a)[r + gamma v(s')]
\end{verbatim}

\section{Ce qu'il faut retenir}

\begin{itemize}[leftmargin=1.6em]
    \item Value Iteration part de l'équation d'optimalité de Bellman.
    \item L'opérateur de Bellman optimal est une contraction pour la norme infinie.
    \item Le théorème du point fixe de Banach garantit l'existence, l'unicité et la convergence vers $v_*$.
    \item Dans FrozenLake, l'algorithme est exact car le modèle de transition est connu.
    \item La politique optimale s'obtient ensuite par maximisation locale dans chaque état.
\end{itemize}

\end{document}
